% META: {"id": "1NSI-WEB-M2", "chapitre": "1NSI-WEB-IHM", "type_objet": "methode", "capacites": ["P-IHM-04A", "P-IHM-04B", "P-IHM-04C"], "capacites_codes": ["C7", "C8", "C9"], "methodes": ["M2"], "mode_creation": "ex_nihilo", "sources_inspiration": [], "status": "generated"}
\begin{fichemethode}{M2}{Lire un formulaire et choisir entre GET et POST}
\textbf{Quand l'utiliser ?} Quand on te montre un formulaire et qu'on te demande ce qui est
transmis, à qui, et comment — ou quand tu dois décider de la méthode à employer dans un
formulaire que tu écris. Une seule lecture répond aux trois questions.

\textbf{Pas à pas :}
\begin{enumerate}
  \item \textbf{Lis la balise \lstinline{<form>}.} Deux attributs seulement comptent :
    \lstinline{action}, qui dit \textbf{à qui} la requête est envoyée, et
    \lstinline{method}, qui dit \textbf{comment}. Sans \lstinline{method}, c'est
    \lstinline{GET} par défaut.
  \item \textbf{Relève chaque champ} et note son attribut \lstinline{name} : c'est ce nom,
    pas l'étiquette affichée, qui voyage jusqu'au serveur. Un champ sans \lstinline{name}
    n'est pas transmis du tout.
  \item \textbf{Déduis la requête.} En \lstinline{GET}, les couples nom-valeur sont collés à
    l'URL après un \lstinline{?}, séparés par des \lstinline{&} :
    \lstinline{/recherche?ville=Tunis&jours=3}. En \lstinline{POST}, ils voyagent dans le
    corps de la requête, invisibles dans la barre d'adresse.
  \item \textbf{Choisis la méthode} avec deux questions. La requête ne fait-elle que
    \textbf{lire} ? \lstinline{GET} : l'URL est partageable, marquable en favori,
    rechargeable sans dommage. Modifie-t-elle quelque chose, ou transporte-t-elle une donnée
    qu'on ne veut pas voir écrite ? \lstinline{POST}.
  \item \textbf{Retiens ce que \lstinline{POST} ne fait pas.} Il n'est pas chiffré. Ce qui
    chiffre, c'est \lstinline{HTTPS}, et lui seul — il protège aussi bien l'URL d'un
    \lstinline{GET} que le corps d'un \lstinline{POST}.
\end{enumerate}

\textbf{Exemple :}

\begin{console}
<form action="/recherche" method="get">
  <label>Ville  <input type="text" name="ville"></label>
  <label>Jours  <input type="number" name="jours"></label>
  <button type="submit">Chercher</button>
</form>
\end{console}

Saisie « Tunis » et « 3 » : le navigateur demande
\lstinline{/recherche?ville=Tunis&jours=3}. Le choix de \lstinline{GET} est correct — la
requête ne fait que consulter, et l'URL obtenue peut être partagée telle quelle.

Le même formulaire avec un champ \lstinline{name="motdepasse"} devrait passer en
\lstinline{POST} : en \lstinline{GET}, le mot de passe s'inscrirait dans la barre d'adresse,
dans l'historique du navigateur et dans les journaux du serveur.

\begin{center}
\begin{tabular}{l|l|l}
& \lstinline{GET} & \lstinline{POST} \\
\hline
où voyagent les données & dans l'URL & dans le corps \\
visible dans l'historique & oui & non \\
URL partageable & oui & non \\
rechargement sans risque & oui & redemande l'envoi \\
chiffré & \multicolumn{2}{c}{seulement si HTTPS} \\
\end{tabular}
\end{center}

\textbf{Pièges :}
\begin{itemize}
  \item \textbf{Croire que \lstinline{POST} chiffre.} Il cache seulement les données de la
    barre d'adresse ; sur le réseau, sans HTTPS, elles circulent en clair tout autant.
  \item Prendre l'étiquette affichée pour le nom transmis. C'est \lstinline{name} qui part,
    et un champ sans \lstinline{name} ne part pas — d'où des valeurs manquantes côté serveur
    sans le moindre message d'erreur.
  \item Employer \lstinline{GET} pour une action qui modifie : le rechargement de la page la
    rejoue, et un lien partagé la déclenche chez quelqu'un d'autre.
  \item Oublier que \lstinline{GET} est la valeur par défaut : un \lstinline{<form>} sans
    \lstinline{method} met tout dans l'URL, y compris ce qu'on croyait discret.
\end{itemize}

\textbf{Vérifier :} soumets le formulaire et lis la barre d'adresse. Si tu vois tes valeurs,
c'est un \lstinline{GET} — et il faut alors se demander si l'une d'elles n'aurait pas dû
rester hors de l'URL.

\textbf{S'entraîner :} exercices C7, C8 et C9 du chapitre.
\end{fichemethode}
